\documentclass[11pt,a4paper]{article}
% Shared preamble for RuPhO English problem statements (match T1 2026 style).
% Use: \input{latex/rupho_en_preamble.tex} after \documentclass.
\usepackage[margin=1in]{geometry}
\usepackage{amsmath,amssymb}
\usepackage{graphicx}
\usepackage{float}
\usepackage{enumitem}
\usepackage{microtype}
\usepackage{caption}
\usepackage{tabularx}
\usepackage{hyperref}

\captionsetup{font=small,labelfont=bf,skip=6pt}

\setlist[itemize]{leftmargin=1.2cm}
\setlist[enumerate]{leftmargin=1.2cm}

% One outer box per subproblem: [id | statement | points] (no inner rules)
\newcommand{\problembox}[3]{%
  \par\addvspace{0.42em}%
  \noindent
  \setlength{\fboxsep}{7.5pt}%
  \setlength{\fboxrule}{0.95pt}%
  \fbox{%
    \begin{minipage}{\dimexpr\linewidth-2\fboxsep-2\fboxrule\relax}
    \begin{tabularx}{\linewidth}{@{}>{\raggedright\arraybackslash\bfseries}p{2.1em} @{\hspace{0.9em}} >{\raggedright\arraybackslash}X @{\hspace{0.9em}} >{\raggedleft\arraybackslash\bfseries}p{2.4em}@{}}
    #1 & #2 & #3
    \end{tabularx}
    \end{minipage}%
  }%
  \par\addvspace{0.32em}%
}


\title{\textbf{Spectroscopy}}
\author{\normalsize Y22 (2022) --- English translation}
\date{}
\begin{document}
\maketitle

Spectroscopic measurements are an important research method that is used in physics, chemistry, biology, medicine and at the intersection of these sciences. Spectral structure of emission, re-emission, absorption, etc. can be studied in different wavelength ranges. For example, using infrared (IR) spectroscopy, one can study the vibrational motion of molecules or their individual fragments or functional groups. Thus, we can talk about the chemical structure of a substance, and, in particular, check the results of chemical synthesis. Visible and near-infrared spectroscopy will be discussed in this problem. You will need to construct a spectrometer and use it to determine the properties of various objects.

\section*{Equipment}

\begin{figure}[H]
    \centering
    \includegraphics[width=0.72\textwidth]{E1_EN_assets/fig1.jpg}
    \caption{Figure 1.}
\end{figure}

Optical bench Movable stands Steel screen Magnetic holders (for lenses and sample holders) Lens with grating Lens without grating RGB LED (three in one housing). Do not exceed 4.5V! Incandescent lamp. Do not exceed 12V! A container for draining used solutions (you can ask to empty them) A container with clean water (for washing the pipette, cuvettes; you can ask to fill it) Magnetic photonic crystal holder Pipette Set of spectrometric cuvettes with lids (12 pcs.), cuvette volume approximately 4 ml Magnetic holder for spectrometric cuvettes Tripod for spectrometric cuvettes (place filled cuvettes in it) Spectroscopic sensor Laptop Power supply Cardboard for making diaphragms Scissors Masking tape (for marks, attaching homemade objects to the installation) Ruler Samples to be tested: S1. Photonic crystal in a Petri dish S2. In 15 ml test tubes: a solution of aniline orange, copper sulfate, solutions with different pH values. In a 2 ml test tube - a solution of thymol blue. S3. A diluted solution of brilliant green. S4. Concentrated solution of brilliant green. Wrapped in film. DO NOT open! Equipment for your convenience: Napkins Table lamp Box for covering the installation if your measurements are disturbed by the light of neighbors

Samples studied:

Equipment for your convenience:

\section*{Spectrometer construction}

The spectrometer will consist of an incandescent lamp (providing a continuous spectrum of radiation), a pair of lenses, a diffraction grating and a one-dimensional array of photodiodes, the data from which is collected into a computer. Attention! When assembling and adjusting the installation, make sure that there are no vertical or horizontal lens shifts: both lenses, the light source and its image (0th maximum) must be on the same axis. Assemble the installation (follow the diagram in Fig. 2): on the voltage source (18), unscrew the current and voltage adjustment knobs to 0; install the incandescent lamp (8) in the tripod (2), apply voltage $12~В$ to the lamp (this is important! the emission spectrum depends on the temperature of the filament = voltage); position the lens without a diffraction grating (6) so that the lamp is in the focus of the lens (the optical power of each lens is +6 diopters); make sure there is a parallel beam of light behind the lens; place one additional stand (2) between the lamp and the lens (it will be used as a sample holder); place the second lens with a diffraction grating (5) at a short distance from the first lens, the diffraction grating should face the first lens; the lines of the diffraction grating must be located vertically; place the screen (3) at the focus of the second lens.

The spectrometer will consist of an incandescent lamp (providing a continuous spectrum of radiation), a pair of lenses, a diffraction grating and a one-dimensional array of photodiodes, the data from which is collected into a computer.

Attention! When assembling and adjusting the installation, make sure that there are no vertical or horizontal lens shifts: both lenses, the light source and its image (0th maximum) must be on the same axis.

Assemble the installation (refer to the diagram in Fig. 2):

\begin{figure}[H]
    \centering
    \includegraphics[width=0.72\textwidth]{E1_EN_assets/fig2.jpg}
    \caption{Figure 2.}
\end{figure}

On the screen at zero maximum you should observe an image of the lamp filament. At the first maximum you should observe the decomposition of the lamp radiation into a spectrum. The vertical size of the spectrum should not be more than 3-4 mm. If the spectrum is not thin, check the installation assembly. Attach the sensor (16) to the screen. The sensor is attached to the screen with magnets. Place the sensitive elements of the sensor (in the middle of the chip) so that all the radiation decomposed into a spectrum (first order) falls on them. The sensor is attached to the screen so that the wires face down. With this configuration, make the red light fall on the right side of the sensor. Because The sensor is quite thick, change the position of the screen slightly so that the sensor (and not the plane of the screen) is in focus of the second lens. In the program running on the computer (17), click the “Get data” button. You should get a graph similar to the one shown below.

\begin{figure}[H]
    \centering
    \includegraphics[width=0.72\textwidth]{E1_EN_assets/fig3.jpg}
    \caption{Figure 3.}
\end{figure}

\section*{Working with the program}

\begin{figure}[H]
    \centering
    \includegraphics[width=0.72\textwidth]{E1_EN_assets/fig4.jpg}
    \caption{Figure 4.}
\end{figure}

Данная програmmа позволяет считывать интенсивность света, падающего на чувствительные элементы датчика (CCD матрицу), а также производить с полученными данными определённые действия. - Основные элементы програmmы : график интенсивности, основные оси, вспомогательная ось для графика отношения, первый график, второй график, график отношения, виджет построения графика, полоса для усреднения, чувствительность датчика, «умный» курсор, виджет для настройки «умного» курсора, кнопки для построения графика отношения первого и второго графиков, блок настройки режима отображения длин волн, кнопка включения/выключения режима отображения длин волн, блок сохранения и загрузки, поле для определения используемого порта - Алгоритм действий для начала работы с програmmой : подключите устройство для считывания к компьютеру через кабель. Выберите на панеле «Используемый порт» (16) – COM3. Снизу под ней должна загореться зелёная надпись «Порт активен». Если этого не произошло, позовите наблюдателя и сообщите о проблеме. - Считывание данных : Для считывания новых данных вам необходимо нажать на кнопку «Получить данные», которая находится в виджете получения данных (7). Существует два слота для хранения данных: первый и второй. В (7) вы можете выбрать в какой именно помещать новые измерения, при этом старые данные, которые хранились в этом слоте безвозвратно исчезнут (если вы их перед этим не сохранили, подробнее об этом в разделе ниже). У вас есть возможность выбрать степень усреднения (8). Число, которое вы выбираете в этом виджете показывает сколько будет сделано измерений перед их усреднением. Усреднение позволяет избежать появление артефактов, а также шумов. Для того, что бы считывать сигналы разной интенсивности вы можете изменять чувствительность (9). При этом вводимое число должно быть делителем числа 50000. - Работа с графиком, График отношения, «Умный» курсор : На координатной сетке (1) одновременно могут отображаться от 0 до 3 графиков (4, 5, 6). 2 графика из слотов для данных (4, 5) и третий (6) – график отношения, являющийся функцией от первых двух: каждая точка этого графика получается делением значения из первого набора данных на второй или наоборот. Так как характерные значения графика отношения на несколько порядков расходятся с основными, добавлена дополнительная ось ординат (3). Построить такой график вы можете с помощью кнопок «Поделить первый график на второй» и «Поделить второй график на первый» (12). При снятии новых данных график отношения пропадает. Для вашего удобства в данную програmmу был добавлен «умный» курсор (10). Этот курсор «ползет» по выбранному графику и отображает координаты текущей точки. Настройка «умного» курсора производится в виджете (11). Там вы можете включить/выключить курсор с помощью одноименной кнопки, а так же выбрать, с первым или вторым графиком будет работать «умный» курсор. При активном графике отношения курсор автоматически наводится на него. - Режим отображения длин волн : в основном режиме по оси X будут отображаться номера пикселей матрицы. Однако с помощью кнопки «Построить зависимость от длин волн» (14) вы можете перестроить ваши данные с учётом откалиброванных длин волн. Для перестройки оси X вам необходимо ввести 3 пары «номер пикселя - длина волны» в соответствующий виджет (13). Построение происходит с помощью метода наименьших квадратов. Если вам нужно перестроить график только по двум опорным длинам волн, то вы можете вместо третьей точки указать значения одной из первых двух или любую их линейную комбинацию. Данные полученной калибровочной зависимости отображаются внизу виджета. - Сохранение и загрузка данных : Для того, что бы сохранить и загрузить, ранее сохранённые данные, вам необходимо использовать виджет (15). При этом с помощью кнопок «Первый» и «Второй» можно выбрать в какой слот вы будете загружать данные или в какой – сохранять. Сохранение происходит в файл с разрешением .dat. Сохранить график отношения нельзя. - Артефакты : иногда на графиках можно увидеть вертикальные полосы. Так бывает когда ваш сигнал очень низкий. Это можно починить, увеличив чувствительность (9) и включив усреднение (8).

- Основные элементы програmmы : график интенсивности, основные оси, вспомогательная ось для графика отношения, первый график, второй график, график отношения, виджет построения графика, полоса для усреднения, чувствительность датчика, «умный» курсор, виджет для настройки «умного» курсора, кнопки для построения графика отношения первого и второго графиков, блок настройки режима отображения длин волн, кнопка включения/выключения режима отображения длин волн, блок сохранения и загрузки, поле для определения используемого порта

- Алгоритм действий для начала работы с програmmой : подключите устройство для считывания к компьютеру через кабель. Выберите на панеле «Используемый порт» (16) – COM3. Снизу под ней должна загореться зелёная надпись «Порт активен». Если этого не произошло, позовите наблюдателя и сообщите о проблеме.

- Считывание данных : Для считывания новых данных вам необходимо нажать на кнопку «Получить данные», которая находится в виджете получения данных (7). Существует два слота для хранения данных: первый и второй. В (7) вы можете выбрать в какой именно помещать новые измерения, при этом старые данные, которые хранились в этом слоте безвозвратно исчезнут (если вы их перед этим не сохранили, подробнее об этом в разделе ниже). У вас есть возможность выбрать степень усреднения (8). Число, которое вы выбираете в этом виджете показывает сколько будет сделано измерений перед их усреднением. Усреднение позволяет избежать появление артефактов, а также шумов. Для того, что бы считывать сигналы разной интенсивности вы можете изменять чувствительность (9). При этом вводимое число должно быть делителем числа 50000.

- Работа с графиком, График отношения, «Умный» курсор : На координатной сетке (1) одновременно могут отображаться от 0 до 3 графиков (4, 5, 6). 2 графика из слотов для данных (4, 5) и третий (6) – график отношения, являющийся функцией от первых двух: каждая точка этого графика получается делением значения из первого набора данных на второй или наоборот. Так как характерные значения графика отношения на несколько порядков расходятся с основными, добавлена дополнительная ось ординат (3). Построить такой график вы можете с помощью кнопок «Поделить первый график на второй» и «Поделить второй график на первый» (12). При снятии новых данных график отношения пропадает. Для вашего удобства в данную програmmу был добавлен «умный» курсор (10). Этот курсор «ползет» по выбранному графику и отображает координаты текущей точки. Настройка «умного» курсора производится в виджете (11). Там вы можете включить/выключить курсор с помощью одноименной кнопки, а так же выбрать, с первым или вторым графиком будет работать «умный» курсор. При активном графике отношения курсор автоматически наводится на него.

- Режим отображения длин волн : в основном режиме по оси X будут отображаться номера пикселей матрицы. Однако с помощью кнопки «Построить зависимость от длин волн» (14) вы можете перестроить ваши данные с учётом откалиброванных длин волн. Для перестройки оси X вам необходимо ввести 3 пары «номер пикселя - длина волны» в соответствующий виджет (13). Построение происходит с помощью метода наименьших квадратов. Если вам нужно перестроить график только по двум опорным длинам волн, то вы можете вместо третьей точки указать значения одной из первых двух или любую их линейную комбинацию. Данные полученной калибровочной зависимости отображаются внизу виджета.

- Сохранение и загрузка данных : Для того, что бы сохранить и загрузить, ранее сохранённые данные, вам необходимо использовать виджет (15). При этом с помощью кнопок «Первый» и «Второй» можно выбрать в какой слот вы будете загружать данные или в какой – сохранять. Сохранение происходит в файл с разрешением .dat. Сохранить график отношения нельзя.

- Артефакты : иногда на графиках можно увидеть вертикальные полосы. Так бывает когда ваш сигнал очень низкий. Это можно починить, увеличив чувствительность (9) и включив усреднение (8).

\section*{Part A: Wavelength Calibration}

Для того, чтобы в дальнейшем можно было работать с количественными величинами длин волн, собранный спектрометр необходимо откалибровать. Т.е. каждому пикселю датчика нужно поставить в соответствие длину волны. В этой задаче известными, опорными значениями будут являться длины волн, на которых интенсивность излучения светодиодов максимальна.

Общий cmысл процедуры калибровки следующий: установить лампу и датчик в рабочее положение, заменить лампу на светодиод и зафиксировать, где находится первый дифракционный порядок излучения светодиода. Таким образом, некоторому пикселю можно будет поставить в соответствие определенную длину волны. Порядок действий при калибровке: установите и включите лампу (8), расположите датчик в первый дифракционный порядок; на экране (3) также должен быть виден нулевой максимум; наклейте малярный скотч (21) на нулевой максимум и отметьте его положение, это будет важно для дальнейшей процедуры калибровки; отключите и снимите лампу, установите вместо нее держатель со светодиодом (7) (светодиод относительно стойки расположен так же как и нить накаливания лампы); на источнике напряжения (18) выкрутите ручки регулировки тока и напряжения в 0; подключите светодиод (7) к блоку питания: красный штекер в «$+$», любой из черных – в «$-$» (в одном корпусе светодиода расположены три излучающих элемента); подайте на светодиод напряжение 3-4 В, он должен загореться (не превышайте напряжение на светодиоде в 4.5 В, сгоревшие светодиоды не заменяются); отрегулируйте положение светодиода по высоте и по углу поворота так, чтобы нулевой максимум совпал с отметкой на малярном скотче, соответствующей нулевому максимуму от лампы; на датчике (16) должно быть видно пятно первого дифракционного порядка; в програmmе, запущенной на компьютере (17), нажмите кнопку «Получить данные»; вы должны увидеть спектр излучения светодиода, похожий на тот, что приведен на рис. 5; если сигнал слишком слабый, увеличьте напряжение на светодиоде (не превышая 4.5 В) или увеличьте чувствительность в програmmе; если сигнал слишком сильный (зашкаливает), уменьшите напряжение на светодиоде или уменьшите чувствительность в програmmе; определите положение максимума (в пикселях) и запишите это число, рядом укажите длину волны соответствующего светодиода (рис. 5); не меняя положения оборудования, переключите светодиод на другой цвет (замените черный штекер на другой), определите положение максимума; излучение разных цветов в светодиоде генерируется немного в разных по высоте местах, и первый дифракционный порядок может немного cmещаться, это нормально, можно не cmещать датчик, если хотя бы часть каждого цветного пятна попадает на него; повторите процедуру для третьего светодиода; в соответствующих полях в програmmе можно ввести пары «номер пикселя-длина волны», и график автоматически будет перестроен в координатах «интенсивность (длина волны)»; можно вернуть на место лампу, проверив ее установку по нулевому максимуму.

Общий cmысл процедуры калибровки следующий: установить лампу и датчик в рабочее положение, заменить лампу на светодиод и зафиксировать, где находится первый дифракционный порядок излучения светодиода. Таким образом, некоторому пикселю можно будет поставить в соответствие определенную длину волны.

Порядок действий при калибровке:

Спектры светодиодов в соответствии с технической документацией. Максимумы излучения: синего --- $\lambda_B = 470~нм$, зеленого --- $\lambda_G = 528~нм$, красного --- $\lambda_R = 627~нм$.

\begin{figure}[H]
    \centering
    \includegraphics[width=0.72\textwidth]{E1_EN_assets/fig5.jpg}
    \caption{Figure 5.}
\end{figure}

\problembox{A1}{Calibrate the spectrometer. Set the lamp voltage to 12 V. Determine the wavelengths corresponding to the minimums shown in Figure 3.}{2.00}

If the positions of key installation parameters (lamp, lenses, screen, sensor) change, the spectrometer, of course, needs to be recalibrated.

Important! In all further work, the samples under study should be placed directly behind the lamp (the gap between the lamp body and the sample is 5-7 mm). It is important that all the light from the lamp that subsequently hits the diffraction grating goes through the sample. Otherwise, the spectrum on the sensor will be a superposition of the spectrum of the object under study and the spectrum of the lamp itself. To further limit the path of unnecessary rays, you can make diaphragms. You are only allowed to attach your own diaphragms to lens bodies. They cannot be attached to cuvettes, sample holders, photonic crystal, etc.

\section*{Radiation transmission study}

In further parts of the problem, the absorption of various objects/solutions/substances at different parameters will be studied. For these objects, the concept of transmittance $t$ can be introduced. It is a function of the wavelength $\lambda$. The analysis of the function $t(\lambda)$ is the basis for the study of the properties and structure of substances. If radiation with wavelength $\lambda$ and intensity $I_0$ is incident on a certain sample, and the intensity of transmitted radiation at the same wavelength is $I_t$ (Fig. 6), then the transmittance coefficient $t$: $$t(\textbackslash\{\}lambda)=\textbackslash\{\}cfrac\{I_t\}\{I_0\}.$$ Технически для нахождения коэффициента пропускания $t(\lambda)$ вам необходимо в програmmе в «первый» график снять спектр исходного излучения $I_0(\lambda)$, затем на пути света расположить образец и во «второй» график снять спектр прошедшего излучения $I_t(\lambda)$. Для вашего удобства в програmmе сделаны кнопки «Поделить первый график на второй» и наоборот. При хорошо снятых данных вы напрямую рассчитаете и cmожете поточечно снять зависимость коэффициента пропускания от номера пикселя (или длины волны, если проведена калибровка). Кнопки деления работают исключительно формально, т.е. делят попиксельно интенсивности и показывают результат (исключая только слишком большие, зашкаливающие значения). Учитывайте это, если ваши данные «шумные» или имеют значительный фон. В таких случаях рекомендуется обработать данные вручную. Конечно, коэффициент пропускания $t$ зависит от ряда параметров: концентрации вещества $n$, толщины слоя вещества $b⟦ M14⟧\varepsilon(\lambda)$ (закон Бугера-Ламберта-Бера): $$t=\textbackslash\{\}exp\{(-n\textbackslash\{\}varepsilon b)\}.$$ Помните, что спектр лампы – не плоский. Помните, что кювета тоже имеет оптические свойства. Поэтому, если вас интересует только поглощение раствора/вещества, то исходным сигналом должен быть свет, прошедший через пустую кювету (или кюветы).

В дальнейших частях задачи будет исследоваться поглощение различных объектов/растворов/веществ при разных параметрах. Для этих объектов можно ввести понятие коэффициента пропускания $t$. Он является функцией длины волны $\lambda$. На анализе функции $t(\lambda)⟦M20 ⟧\lambda$ и интенсивностью $I_0$, а интенсивность прошедшего излучения на той же длине волны равна $I_t$ (рис. 6.), то коэффициент пропускания $t$: $$t(\textbackslash\{\}lambda)=\textbackslash\{\}cfrac\{I_t\}\{I_0\}.$$ Технически для нахождения коэффициента пропускания $t(\lambda)$ вам необходимо в програmmе в «первый» график снять спектр исходного излучения $I_0(\lambda )$, затем на пути света расположить образец и во «второй» график снять спектр прошедшего излучения $I_t(\lambda)$. Для вашего удобства в програmmе сделаны кнопки «Поделить первый график на второй» и наоборот. При хорошо снятых данных вы напрямую рассчитаете и cmожете поточечно снять зависимость коэффициента пропускания от номера пикселя (или длины волны, если проведена калибровка). Кнопки деления работают исключительно формально, т.е. делят попиксельно интенсивности и показывают результат (исключая только слишком большие, зашкаливающие значения). Учитывайте это, если ваши данные «шумные» или имеют значительный фон. В таких случаях рекомендуется обработать данные вручную. Конечно, коэффициент пропускания $t$ зависит от ряда параметров: концентрации вещества $n$, толщины слоя вещества $b$, через который прошел свет, и некоторой характеристической величины – коэффициента экстинкции $\varepsilon(\lambda)$ (закон Бугера-Ламберта-Бера): $$t=\textbackslash\{\}exp\{(-n\textbackslash\{\}varepsilon b)\}.$$ Remember that the spectrum of a lamp is not flat. Remember that the cuvette also has optical properties. Therefore, if you are only interested in the absorption of a solution/substance, then the initial signal should be the light passing through an empty cuvette (or cuvettes).

\begin{figure}[H]
    \centering
    \includegraphics[width=0.72\textwidth]{E1_EN_assets/fig6.jpg}
    \caption{Figure 6.}
\end{figure}

\section*{Part B: Photonic Crystal (or APhO-17 in 20 minutes)}

In this part of the task we will examine the photonic crystal given to you. A photonic crystal is a structure in which the refractive index of a material changes periodically on scales on the order of the wavelength of light. Because of this structure, photonic crystals at certain wavelengths practically stop transmitting light and only reflect it. It is known from wave optics that if the structure of such a crystal in the direction of incidence of light (angle of incidence $\theta = 0$) has a spatial period $D$, then a noticeable change in transmission can be observed only at wavelengths $\lambda$, determined by the equality:\textbackslash\{\}[2Dn=m\textbackslash\{\}lambda\textbackslash\{\}] - the so-called Bragg-Snell law, where $m$ is the natural number, and $n$ is the refractive index of the material (dispersion can be neglected).

\problembox{B0}{The photonic crystal you were given (see equipment (S1)) is very fragile. Insert the magnetic holder (4) into the stand (2). Install the photonic crystal holder (11) onto the magnetic holder (4). Remove the photonic crystal (S1) from the Petri dish and place it on the holder (11). Do not drop or poke the crystal. Damage to the crystal will result in disqualification for this round.}{-20.00}

The photonic crystal given to you is manufactured in such a way that it has 4 deep transmission minima in the visible range. (Kirill S. Napolskii, Alexey A. Noyan, Sergey E. Kushnir. Control of high-order photonic band gaps in one-dimensional anodic aluminum photonic crystals // Optical Materials 109 (2020) 110317)

\problembox{B1}{Take the spectrum of the original radiation. Place the photonic crystal in the holder along the beam path in the system (Fig. 2). Take the spectrum of transmitted radiation. Observe four deep transmittance minima. Determine at what wavelengths $\lambda$ the transmission minima are observed. Determine which values ​​of $m$ they correspond to. Determine the corresponding values ​​of $t$.}{1.60}

It follows from the Bragg-Snell law that intensity minima should also exist at larger values ​​of $m$.

\problembox{B2}{Find another minimum of photonic crystal transmission. Which $\lambda$, $m$ and $t$ correspond to it? Save the graph used to solve this item in the folder Plots. “Last name” with the name “B2”. The chart must contain the found minimum and at least one other one obtained earlier. If there is no file, the item will not be evaluated!}{0.40}

Similarly, transmission minima should also be observed at lower $m$. It turns out that a total of 7 transmission minima can be detected in the region of wavelengths emitted by the lamp.

\problembox{B3}{Find the sixth minimum. Describe how you do this. Which $\lambda$, $m$ and $t$ correspond to it? Save the graph used to solve this item in the Plots_“Last Name” folder with the name “B3”. The chart must contain the found minimum and at least one other one obtained earlier. If the file is missing, the item will not be evaluated.}{0.50}

\problembox{B4}{Find the seventh minimum. Describe how you do this. Find the corresponding $\lambda$ and $m$. Find the transmittance $t$. Save the graph used to solve this item in the Plots_“Last Name” folder with the name “B4”. The chart must contain the found minimum and at least one other one obtained earlier. If the file is missing, the item will not be evaluated.}{1.00}

\problembox{B5}{Use the Bragg-Snell law to find the value of $Dn$ for the photonic crystal given to you.}{0.50}

\section*{Part C. Band-pass optical filter (or how much IPhO-05 lies)}

When studying various samples, the need arises to create narrow-spectrum radiation. Wavelength tunable lasers are very expensive, and when there is no need for monochromatic and coherent radiation, a combination of filters can be used. Filters will highlight a narrow range from a continuous spectrum. In this part of the problem, two substances (solutions of aniline orange and copper sulfate) will act as a short-wave and long-wave filter. And their mixture is stripe.

For this part of the work you will need solutions of aniline orange and copper sulfate (S2). They will need to be poured into cuvettes (13) using a pipette (12). If you need to pipette another solution, rinse it with water before doing so (10). Cover the cuvettes with lids to prevent liquid from spilling. For spectroscopic examination of solutions, insert the cuvette(s) into the holder (14) and attach to the magnetic holder (4) in the rack (2) (Fig. 7). If there is little solution in the cuvette or the source light affects the measurements, consider creating limiting diaphragms from cardboard (19). The cuvettes can also be washed. Drain dirty water into container (9). Similar steps must be performed in parts D and E.

For this part of the work you will need solutions of aniline orange and copper sulfate (S2). They will need to be poured into cuvettes (13) using a pipette (12). If you need to pipette another solution, rinse it with water before doing so (10). Cover the cuvettes with lids to prevent liquid from spilling. For spectroscopic examination of solutions, insert the cuvette(s) into the holder (14) and attach to the magnetic holder (4) in the rack (2) (Fig. 7). If there is little solution in the cuvette or the source light affects the measurements, consider creating limiting diaphragms from cardboard (19). The cuvettes can also be washed. Drain dirty water into container (9). Similar steps must be performed in parts D and E.

\begin{figure}[H]
    \centering
    \includegraphics[width=0.72\textwidth]{E1_EN_assets/fig7.jpg}
    \caption{Figure 7.}
\end{figure}

\problembox{C1}{Fill the spectrometric cuvette with aniline orange solution (approximately 3 ml). Remove the dependence of the transmittance of aniline orange $t_a(\lambda)$. Plot the spectroscopic template < > on your answer sheet. The wavelength and transmittance scales are plotted for you. Indicate whether aniline orange is a short-wave filter (passes short waves) or long-wave (passes long waves).}{0.50}

\problembox{C2}{Fill the spectrometric cuvette with copper sulfate solution (approximately 3 ml). Remove the dependence of the transmittance of copper sulfate $t_{Cu}(\lambda)$. Plot the same spectroscopic template < > on the answer sheet. The wavelength and transmittance scales are plotted for you. Indicate what type of filter the copper sulfate solution is: short-wave (passes short waves) or long-wave (passes long waves).}{0.50}

\problembox{C3}{Fill the spectrometric cuvette with a mixture of solutions of aniline orange and copper sulfate in a 1:1 ratio (i.e., 1.5 ml of each). Remove the dependence of the mixture transmittance $t_f(\lambda)$. Plot the spectroscopic pattern "SpC3" on the answer sheet. The wavelength and transmittance scales are plotted for you. Save the graph used to solve this item in the folder Plots_“Last Name” with the name “C3”. If there is no file, the item will not be evaluated!}{0.50}

\problembox{C4}{Find theoretically how $t_a, t_{Cu}$ and $t_f$ are related.}{1.00}

\problembox{C5}{Calculate the values ​​of $t_f^{теор}(\lambda)$ using the formula you suggested from point C4.  Plot $t_f^{теор}(\lambda)$ on the same spectroscopic template < > on the answer sheet. Draw a conclusion about the validity of your formula from point C4. Draw a conclusion: is it possible to obtain a bandpass filter using short- and long-wave filters?}{1.50}

\section*{Part D. What color is the green stuff? (or button accordion from Kvant)}

Read carefully again how to handle solutions and cuvettes at the beginning of Part C.

\problembox{D1}{Take a test tube $Z$ (see equipment S4, film-wrapped tube) containing a solution of tetraethyl 4,4-diaminotriphenylmethane oxalate. Depending on the tilt, a layer of liquid of varying thickness can be created. Remove the lamp from the tripod (carefully! the lamp is hot!) and place it horizontally on the bench. Look at the burning lamp through a layer of liquid of maximum thickness. Specify the observed color. Look at the burning lamp through a thin layer of liquid. Specify color. How does color change depending on thickness?}{0.50}

\problembox{D2}{Using a diluted solution of brilliant green (see equipment S3), record the transmission spectra of brilliant green depending on the thickness of the liquid layer. The concentration of the diluted solution is 10 mg/l, which is 1000 times less than in test tube $Z$. Plot these transmittance spectra on one graph. Save the graphs used to solve this item in the folder Plots_“Last Name” with the name “D3_N”, where N is the thickness of the liquid layer in millimeters. If there are no files, the item will not be evaluated!}{2.00}

\problembox{D3}{For two wavelengths ($\lambda_1=480~нм$ and $\lambda_2=790~нм$), write in the table the dependence of the transmittance on the path traveled by the light. Determine the extinction coefficients $\varepsilon_1$ and $\varepsilon_2$ for these two wavelengths.}{2.00}

\problembox{D4}{Explain the effect you observed in D1.}{0.50}

\section*{Part E. pH indicator (we’ll give this to IJSO students later)}

The pH value (pH, minus the logarithm of the proton concentration in solution) is the most important parameter of various chemical and biological fluids. To determine it, special indicator substances are used, the transmission spectrum of which is sensitive to the concentration of protons. One such indicator is thymol blue.

In this part of the problem, it is necessary to remove the dependence of transmission spectra on the pH of solutions. Fill a cuvette with 3 ml of a solution with a certain pH value (see equipment S2). Using a pipette, add two (2) drops of thymol blue (S2) to the cuvette. Pour excess thymol back into the test tube or waste container. Rinse the pipette with water (10). Immerse the pipette in the cuvette with a mixture of solution and thymol. Pipette the liquid in and out and stir the solution. The solution is ready for spectroscopic measurements. When moving to another solution, try to take the same amount and drop the same drops of thymol in the same volume, i.e. try to maintain a constant concentration of thymol. Otherwise, some normalization will be required (if necessary, this can be done based on data in the IR range). Read carefully again how to handle solutions and cuvettes at the beginning of Part C.

\problembox{E1}{Take transmission spectra for 8 pH values ​​of solutions. Plot all these spectra on one spectroscopic template < >.}{1.60}

\problembox{E2}{Indicate the wavelength $\lambda_{ch}$ in the visible range at which thymol blue is most sensitive to pH. Indicate the wavelength $\lambda_{iso}$ in the visible range at which thymol blue is least sensitive to pH.}{0.60}

\problembox{E3}{For the wavelength at which thymol blue is most sensitive to pH, plot the extinction coefficient versus pH. Express the extinction coefficient in conventional units.}{2.00}

\problembox{E4}{Suggest a way to spectrally determine the pH value of a solution to which a certain amount of thymol blue has been added. What is the minimum set of data you might need?}{0.80}

\vspace{1em}
\noindent\textit{Source:} \url{https://pho.rs/p/2847}

\end{document}